%% chapters/chapter03/chapter03.tex
%% CHAPTER 3: METHODOLOGY

\chapterfigpath{chapter03}

\chapter{Methodology}
\label{chap:methodology}

\section{Research Design}
\label{sec:research_design}

This research adopts an engineering-design methodology to develop, validate, and deploy HippoCortex, a biologically-inspired continual learning framework. The research is structured in two major stages: Stage 1 focuses on designing, implementing, and validating the core CL algorithms on sequence benchmarks in simulation; Stage 2 extends this architecture to a hybrid model for real-time, raw-data-free deployment on a physical quadruped robot.

\subsection{Research Approach and Workflow}
The workflow is divided into four iterative phases to align theoretical modeling with codebase validation:
\begin{enumerate}
    \item \textbf{Analysis Phase:} Reviewing memory and compute constraints on edge robotic devices and defining mathematical requirements.
    \item \textbf{Design Phase:} Formulating the SWR stats buffer, CVAE architecture, null-space gradient projection equations, and hybrid Mamba-Transformer layers.
    \item \textbf{Implementation Phase:} Constructing the Python modules, training trainer routines, and configuring on-board middleware.
    \item \textbf{Evaluation Phase:} Validating the model on Split-CIFAR100 and ImageNet-R \cite{hendrycksManyFacesRobustness2021}, and benchmark testing on-device latency and accuracy on the Unitree Go2 Edu.
\end{enumerate}

Figure~\ref{fig:development_workflow} illustrates the training and consolidation workflow.

\begin{figure}[ht]
    \centering
    \resizebox{0.95\textwidth}{!}{%
\begin{tikzpicture}[
    node distance=1.2cm and 1.0cm,
    box/.style={
        draw, 
        rectangle, 
        rounded corners=5pt, 
        minimum width=2.8cm, 
        minimum height=0.9cm, 
        align=center, 
        thick,
        font=\sffamily\small
    },
    phasebox/.style={box, fill=blue!5, draw=blue!70!black},
    stepbox/.style={box, fill=gray!2, draw=gray!40},
    evalbox/.style={box, fill=green!5, draw=green!70!black},
    arrow/.style={-Stealth, thick, draw=gray!80},
    labelnode/.style={font=\sffamily\scriptsize, color=gray!90!black}
]

% Nodes
\node[stepbox] (start) {Task stream input\\$D_t$};
\node[phasebox, right=of start] (warmup) {Warmup Phase\\(Train head on $D_t$)};
\node[phasebox, right=of warmup] (joint) {Joint Phase\\(Train with replay)};
\node[phasebox, right=of joint] (consolidation) {Sleep Phase\\(Consolidation)};
\node[evalbox, below=of consolidation] (eval) {Evaluation\\(AA, AIA, FM)};

% Connect
\draw[arrow] (start) -- (warmup);
\draw[arrow] (warmup) -- (joint);
\draw[arrow] (joint) -- (consolidation);
\draw[arrow] (consolidation) -- (eval);

% Feedback loop to next task
\draw[arrow] (eval.west) -| (start.south)
    node[pos=0.25, above, labelnode] {Next task $t+1$};

\end{tikzpicture}%
}

    \caption{Development and training workflow of the HippoCortex framework, detailing the transition between awake learning phases and the sleep consolidation phase.}
    \label{fig:development_workflow}
\end{figure}

\section{Development and Deployment Environment}
\label{sec:development_environment}

To ground the empirical findings, development and testing are split between a cloud GPU instance for training the core sequence models and an on-board edge computer for real-world quadruped deployment.

\subsection{Hardware Specifications}
Table~\ref{tab:hardware} lists the hardware configurations of the experimental systems. The NVIDIA A6000 cloud instance provides the VRAM and compute capacity needed to run the 16-hour Mamba-CL baseline training and evaluate HippoCortex models in simulation. The Jetson Orin NX (16GB) on the Unitree Go2 Edu runs inference on-board, without relying on off-site servers.

\begin{table}[ht]
    \centering
    \caption{Hardware specifications of the experimental and deployment platforms.}
    \label{tab:hardware}
    \resizebox{\textwidth}{!}{%
    \begin{tabular}{lll}
        \toprule
        \textbf{Component} & \textbf{Cloud Training Instance (Stage 1)} & \textbf{Robotic Edge Platform (Stage 2)} \\
        \midrule
        Host System & Thunder Compute Cloud Instance & Unitree Go2 Edu Quadruped \\
        Processor & 4 vCPUs & Jetson Orin NX (16GB RAM, 100 TOPS AI) \\
        GPU & NVIDIA A6000 (48GB VRAM) & Integrated NVIDIA Ampere GPU \\
        System RAM & 32 GB & 16 GB LPDDR5 (shared) \\
        Storage & 100 GB SSD & 128 GB M.2 NVMe SSD \\
        Sensors & N/A & L1 Lidar, RealSense D435i RGB-D, IMU \\
        Operating System & Ubuntu 22.04 LTS & Ubuntu 20.04 LTS with ROS2 Foxy \\
        \bottomrule
    \end{tabular}%
    }
\end{table}

\subsection{Software Tools and Libraries}
The software stack is built using:
\begin{itemize}
    \item \textbf{Primary Language:} Python 3.10, ensuring compatibility with PyTorch, CUDA, and robotics middleware.
    \item \textbf{Deep Learning Backend:} PyTorch 2.0+ (compiled with CUDA 11.8), leveraging its optimized \code{torch.linalg} backend for fast SVD and QR decomposition operations during null-space projector updates.
    \item \textbf{Sequence Backbone:} The official \code{mamba-ssm} library, utilizing custom CUDA kernels for selective scan operations.
    \item \textbf{Robot Middleware:} ROS2 Foxy Fitzroy and CycloneDDS for on-board sensor ingestion and command routing.
    \item \textbf{Robotics SDK:} Unitree SDK2, enabling the Python controller to send joint and velocity commands directly to the quadruped motor boards.
\end{itemize}

\section{System Architecture}
\label{sec:system_architecture}

Figure~\ref{fig:system_architecture} illustrates the system architecture of HippoCortex, highlighting the on-board components running on the Jetson Orin NX edge platform.

\begin{figure}[ht]
    \centering
    \resizebox{0.95\textwidth}{!}{%
\begin{tikzpicture}[
    node distance=1.4cm and 1.2cm,
    box/.style={
        draw, 
        rectangle, 
        rounded corners=5pt, 
        minimum width=3.0cm, 
        minimum height=1.0cm, 
        align=center, 
        thick,
        font=\sffamily\small
    },
    sensorbox/.style={box, fill=orange!5, draw=orange!70!black},
    backbonebox/.style={box, fill=blue!5, draw=blue!70!black},
    hcbox/.style={box, fill=green!5, draw=green!70!black},
    controlbox/.style={box, fill=red!5, draw=red!70!black},
    arrow/.style={-Stealth, thick, draw=gray!80},
    labelnode/.style={font=\sffamily\scriptsize, color=gray!90!black}
]

% Nodes
\node[sensorbox] (sensors) {Sensors (Go2 Edu)\\RGB-D, LiDAR, IMU};
\node[backbonebox, right=of sensors, xshift=0.5cm] (mamba) {Mamba Block\\(Temporal State)};
\node[backbonebox, below=of mamba] (transformer) {Transformer Block\\(Spatial Attention)};
\node[hcbox, right=of mamba, xshift=0.5cm] (swr) {SWR Generator\\(CVAE Replay)};
\node[hcbox, below=of swr] (projector) {Null-Space Projector\\(Gradient Orthogonalizer)};
\node[controlbox, right=of swr, xshift=0.5cm] (ros) {ROS2 SDK\\CycloneDDS};
\node[controlbox, below=of ros] (motors) {Joint Control\\Unitree SDK2};

% Connect Sensors to Backbone
\draw[arrow] (sensors.east) -- (mamba.west);
\draw[arrow] (sensors.east) |- (transformer.west);

% Connect Backbone layers
\draw[arrow] (mamba.south) -- (transformer.north);
\draw[arrow] (transformer.east) -- (projector.west);
\draw[arrow] (mamba.east) -- (swr.west);

% Connect HippoCortex CL components to Control/ROS
\draw[arrow] (swr.east) -- (ros.west);
\draw[arrow] (projector.east) -- (motors.west);
\draw[arrow] (ros.south) -- (motors.north);

% Frame to show boundaries (Edge platform)
\draw[dashed, draw=gray!60, thick] (2.3, 0.8) rectangle (11.75, -3.8);
\node[labelnode, anchor=south] at (7.0, -3.7) {\textbf{Jetson Orin NX (On-Board Edge Deployment)}};

\end{tikzpicture}%
}

    \caption{System architecture of HippoCortex showing the flow from physical sensor ingestion on the Unitree Go2 Edu to the hybrid sequence backbone, supervised by the stats buffer, SWR generator, and null-space projector.}
    \label{fig:system_architecture}
\end{figure}

\subsection{Mamba Sequence Backbone}
The sequence backbone processes sequential patch embeddings or sensor data. The network consists of an $L$-layer stack of selective state-space blocks (detailed in the codebase file \code{hippocortex/}\allowbreak\code{models/}\allowbreak\code{backbone.py}). 

Currently, the primary backbone is configured using the \textbf{De-focus Mamba} \cite{taoLearning1DCausal} architecture. De-focus Mamba treats images as a 1D causal sequence of tokens, utilizing learnable bandpass filters to broaden the attention range. This design mitigates the ``over-focus'' issue inherent to applying standard causal language Mamba directly to visual data, where the network gradients and activation weights tend to concentrate on a very small fraction of visual tokens, limiting the representation capacity. 

For future experimental stages, the framework is designed to also support \textbf{Vision Mamba (Vim)} \cite{zhuVisionMambaEfficient2024}. Unlike the 1D causal scanning of De-focus Mamba, Vim uses bidirectional Mamba blocks to scan patch sequences in both forward and backward directions, combined with position embeddings to capture non-causal 2D global context. In future experiments, we will compare 1D causal scanning against bidirectional 2D scanning to analyze how different state scan directions affect the stability-plasticity trade-off and representation drift under null-space projection.

For both models, a forward hook is registered on the penultimate block to capture the hidden state activations. The raw captured tensor has shape $(B, T_s, d_{\text{model}})$, where $B$ is the batch size, $T_s$ is the sequence length, and $d_{\text{model}}$ is the model dimension. To obtain a sequence-level summary, the hook mean-pools the activations across the temporal dimension $T_s$ to yield the representation $H \in \mathbb{R}^{B \times d_{\text{model}}}$:
\begin{equation}
    H = \frac{1}{T_s} \sum_{k=1}^{T_s} h_k^L
    \label{eq:hidden_extraction}
\end{equation}
where $T_s$ is the sequence length, and $h_k^L \in \mathbb{R}^{B \times d_{\text{model}}}$ is the output of the final Mamba block at step $k$. This representation $H$ is fed to the classification head during the awake phase and serves as the training target for consolidation during the sleep phase.

\subsection{Stats Buffer}
The stats buffer (detailed in \code{stats\_buffer.py}) acts as a compact, raw-data-free long-term memory. At the end of each task $t$, the real hidden states $H_t$ are collected. Instead of saving these states, the buffer stores their mean $\mu_t$ and variance $\sigma_t^2$ (both in $\mathbb{R}^{d_{\text{model}}}$):
\begin{equation}
    \mu_{t} = \frac{1}{N} \sum_{i} H_{t,i}, \ \sigma_{t}^2 = \frac{1}{N} \sum_{i} (H_{t,i} - \mu_{t})^2
    \label{eq:buffer_stats}
\end{equation}
Because the buffer stores only two $d_{\text{model}}$-dimensional vectors per task, the memory footprint scales as $\mathcal{O}(T \times d_{\text{model}})$. For $T=20$ tasks and $d_{\text{model}}=128$, the footprint is less than $20$\,KB, making it highly suitable for resource-constrained edge systems.

\subsection{Sharp Wave Ripple (SWR) CVAE Generator}
The SWR generator (detailed in \code{hippocortex/}\allowbreak\code{models/}\allowbreak\code{swr\_}\allowbreak\code{generator.py}) is structured as a Conditional Variational Autoencoder (CVAE). The encoder maps concatenated hidden states and task embeddings (dimension $d_{\text{model}} + d_{\text{task\_emb}}$) to latent parameters $(\mu_q, \log\sigma_q^2)$ using two 512-unit layers. The decoder decodes latent samples $z \sim \mathcal{N}(\mu_q, \sigma_q^2)$ back to the $d_{\text{model}}$-dimensional space.

The model is trained by optimizing the Evidence Lower Bound (ELBO) loss \cite{sohnLearningStructuredOutput2015}:
\begin{equation}
    \mathcal{L}_{\text{ELBO}} = \|H - \hat{H}\|^2 + \beta D_{\text{KL}}\left(\mathcal{N}(\mu_q, \sigma_q^2) \parallel \mathcal{N}(0, I)\right)
    \label{eq:elbo_loss}
\end{equation}
where $\hat{H}$ is the reconstructed state and $\beta = 1.0$. During the consolidation phase, the generator draws samples $z \sim \mathcal{N}(\mu_t, \sigma_t^2)$ from the statistics stored in the stats buffer and decodes them to generate synthetic hidden states $\hat{H}_{\text{past}}$, simulating hippocampal sharp-wave ripples. Note the operational difference: during active training, the encoder computes the posterior distribution $q(z \mid H, c)$ dynamically to output the latent distribution $(\mu_q, \sigma_q^2)$, which is regularized toward the standard normal prior $\mathcal{N}(0, I)$ by the KL divergence. During offline sleep replay, however, the encoder is bypassed entirely. The latent vector $z$ is sampled directly from the Gaussian distribution $\mathcal{N}(\mu_t, \sigma_t^2)$ defined by the statistics buffer for task $t$, and passed to the decoder to reconstruct the synthetic hidden states.

\subsection{Null-Space Gradient Projector}
The null-space gradient projector (detailed in \code{null\_space\_projector.py}) constrains parameter updates to prevent interference. Given a batch of hidden states, the projector performs SVD to extract the primary feature directions \cite{sahaGradientProjectionMemory2021}:
\begin{equation}
    H_t = U_h \Sigma V_h^\top
    \label{eq:svd_states}
\end{equation}
We keep the singular vectors in $V_h$ that capture $99\%$ of the variance, merging them with the prior basis $U_{t-1}$ via QR decomposition to form the new basis $U_t$ \cite{luVisualPromptTuning}:
\begin{equation}
    Q, R = \text{QR}\left(\left[U_{t-1}, V_{h, 1:r}\right]\right)
    \label{eq:qr_merge}
\end{equation}
The projector truncates $Q$ to a maximum size defined by the $\text{rank\_budget}$ ($k \le 200$) to yield $U_t$. During joint training, gradients $g$ are projected into the null-space \cite{sahaGradientProjectionMemory2021}:
\begin{equation}
    g_{\text{proj}} = g - (g U_t) U_t^\top
    \label{eq:projection_step}
\end{equation}
which is applied to Mamba's time-invariant parameter matrices---specifically the input projection $W_{\text{in}}$, the output projection $W_{\text{out}}$, and the discretized state transition matrix $\bar{A}$---where features are linearly mapped, preserving prior task representations.

\section{Codebase Structural Validation}
\label{sec:codebase_validation}

To ensure the empirical results are reproducible, the methodology is verified directly against the python codebase implementation. Table~\ref{tab:codebase_validation} provides a verification matrix mapping each component to its codebase file, class, and hyperparameter configuration.

\begin{table}[ht]
    \centering
    \caption{Codebase verification matrix validating that the theoretical methodology maps directly to the implemented classes, parameters, and shapes in the Python codebase.}
    \label{tab:codebase_validation}
    \resizebox{\textwidth}{!}{%
    \begin{tabular}{llll}
        \toprule
        \textbf{Component} & \textbf{Theory / Formulation} & \textbf{Codebase File} & \textbf{Verified Specifications \& Hyperparameters} \\
        \midrule
        Sequence Backbone & Stacked Selective SSM & \code{hippocortex/}\allowbreak\code{models/}\allowbreak\code{backbone.py} & \code{MambaBackbone} with $d_{\text{model}}=128$, $n_{\text{layers}}=48$, $d_{\text{state}}=16$ \\
        Stats Buffer & $\mathcal{O}(T \times d_{\text{model}})$ Mean/Var Store & \code{hippocortex/}\allowbreak\code{cl/}\allowbreak\code{stats\_buffer.py} & \code{StatsBuffer} storing $(\mu_t, \sigma_t^2) \in \mathbb{R}^{d_{\text{model}}}$ \\
        SWR Generator & Conditional VAE (CVAE) & \code{hippocortex/}\allowbreak\code{models/}\allowbreak\code{swr\_generator.py} & \code{SWRGenerator} with latent $d_{\text{latent}}=64$, MLP hidden $[512, 512]$ \\
        Null-Space Projector & Orthogonal Projector $P = I - U U^\top$ & \code{hippocortex/}\allowbreak\code{cl/}\allowbreak\code{null\_space\_projector.py} & \code{NullSpaceProjector} with $\text{rank\_budget}=200$, threshold $\ge 99\%$ var \\
        Consolidation Driver & 3-phase Sleep-Wake Replay & \code{hippocortex/}\allowbreak\code{cl/}\allowbreak\code{consolidation.py} & \code{consolidate()} routine using 512 synthetic samples \\
        \bottomrule
    \end{tabular}%
    }
\end{table}


\section{Implementation Details}
\label{sec:implementation}

The three-phase training loop is coordinated by the \code{Trainer} class in \code{trainer.py}. Listing~\ref{lst:trainer_loop} shows this Python routine.

\begin{lstlisting}[language=Python, caption={Continual training coordination routine from \code{hippocortex/}\allowbreak\code{training/}\allowbreak\code{trainer.py}.}, label={lst:trainer_loop}]
def train_task(self, task_id: int, loader: DataLoader) -> dict[str, float]:
    n_classes = self.cfg.dataset.n_classes_per_task

    # Phase 1: Awake - Warmup
    self.backbone.set_task_head(n_classes)
    self.backbone.train()
    warmup_loss = self._run_warmup(loader)

    # Phase 2: Awake - Joint Training
    joint_loss = 0.0
    if task_id > 0:
        joint_loss = self._run_joint_phase(task_id, loader)

    # Phase 3: Sleep - Consolidation
    current_hidden = self._extract_all_hidden(loader)
    replay_loss = consolidate(
        self.backbone, self.swr_gen, self.projector, self.buffer,
        task_id, self.optimizer, current_hidden,
        n_samples=self.cfg.training.n_replay_samples
    )

    return {
        "task_id": float(task_id),
        "loss_warmup": warmup_loss,
        "loss_joint": joint_loss,
        "replay_loss": replay_loss
    }
\end{lstlisting}

\section{Data Collection and Preprocessing}
\label{sec:data_collection}

\subsection{Data Sources}
The Stage 1 core algorithm is evaluated on two standard image sequence datasets:
\begin{itemize}
    \item \textbf{Split-CIFAR100:} Composed of $20$ sequential tasks, with each task containing $5$ classes from the CIFAR-100 dataset. This benchmark tests the model's capacity to handle class-incremental drift over long task streams.
    \item \textbf{ImageNet-R:} Contains $16$ tasks of style and rendition shifts (paintings, cartoons, embroidery) of ImageNet classes \cite{hendrycksManyFacesRobustness2021}, testing the model's robustness under domain and distribution shift.
\end{itemize}

\subsection{Data Preprocessing}
Images are preprocessed before being fed to the Mamba backbone:
\begin{enumerate}
    \item \textbf{Resolution Standardisation:} Images are resized to $224 \times 224$ pixels.
    \item \textbf{Patch Embedding:} Following standard vision state-space protocols \cite{taoLearning1DCausal, zhuVisionMambaEfficient2024}, the resized images are split into non-overlapping patches of size $16 \times 16$ pixels, which are projected into a 1D sequence of length $196$ with embedding dimension $d_{\text{model}}=128$.
    \item \textbf{Normalization:} Patch vectors are normalized using the mean and standard deviation computed over the ImageNet training set.
\end{enumerate}

\section{Testing and Evaluation Strategy}
\label{sec:testing_strategy}

\subsection{Evaluation Metrics}
The framework is evaluated using three continual learning metrics (detailed in Section 2.5): Average Accuracy ($AA$), Average Incremental Accuracy ($AIA$), and Forgetting Measure ($FM$). 

On-device edge validation measures:
\begin{itemize}
    \item \textbf{Inference Latency:} Measured in milliseconds per forward pass on the Jetson Orin NX.
    \item \textbf{Memory Footprint:} Monitored in Megabytes (MB) of VRAM and system RAM during training and evaluation phases.
\end{itemize}

\subsection{Testing Approach}
The testing strategy combines simulated validation and on-device robotic testing:
\begin{itemize}
    \item \textbf{Unit and Integration Testing:} The stats buffer and null-space projection operations are validated using automated unit tests. These are implemented in \code{tests/}\allowbreak\code{test\_}\allowbreak\code{stats\_}\allowbreak\code{buffer.py} and \code{tests/}\allowbreak\code{test\_}\allowbreak\code{projector.py} to verify mathematical correctness.
    \item \textbf{On-Device Robotic Testing:} In Stage 2, the hybrid model is validated on the Go2 Edu platform in a multi-task terrain and object recognition task stream.
\end{itemize}

\section{Chapter Summary}
\label{sec:methodology_summary}

This methodology establishes a structured pipeline connecting our continuous-time mathematical formulations with the actual software architecture of the Python modules. By defining a clear mapping between the theoretical components---such as the stats buffer, CVAE, and gradient orthogonalization---and their concrete Python implementations, we ensure both mathematical verification and hardware execution profiling are grounded. With the software stack and experimental framework thus validated, we proceed in Chapter 4 to analyze the empirical behavior of our baseline reproduction under these design constraints.
